\subsection{Księga pierwsza -- twierdzenia}
\setcounter{counter_euclid}{0}

% Twierdzenie 1.
\begin{euclid_theorem}
    Na danej linii prostej skonstruuj trójkąt równoboczny o żadanych bokach.
\end{euclid_theorem}

% Twierdzenie 2.
\begin{euclid_theorem}
    Skonstruuj odcinek równy danemu odcinkowi którego koniec jest zadanym punktem.
\end{euclid_theorem}

% Twierdzenie 3.
\begin{euclid_theorem}
    Mając dane dwie linie proste nierówne, od większej odciąć linię równą mniejszej.
\end{euclid_theorem}

% Twierdzenie 4.
\begin{euclid_theorem}
    Jeśli dwa trójkąty mają dwa boki odpowiednio równe dwóm innym, i jeżeli kąty zawarte między bokami równoległymi są równe, wtedy ich podstawy również są sobie równe i pozostałe kąty równe są odpowiednim kątom.
\end{euclid_theorem}

% Twierdzenie 5.
\begin{euclid_theorem}
    W trójkątach równoramiennych kąty przy podstawie są sobie równe oraz kąty powstałe przez przedłużenie boków równych są sobie równe.
\end{euclid_theorem}

% Twierdzenie 6.
\begin{euclid_theorem}
    Jeżeli w trójkącie dwa kąty są sobie równe, to boki leżące naprzeciw równych kątów są sobie równe.
\end{euclid_theorem}

% Twierdzenie 7.
\begin{euclid_theorem}
    Na tej samej podstawie i z tej samej strony nie mogą być wykreślone dwa trójkąty takie, żeby boki w tych trójkątach przy obydwu końcach wspólnej podstawy były między sobą równe.
\end{euclid_theorem}

% Twierdzenie 8.
\begin{euclid_theorem}
    Jeżeli dwa boki jednego trójkąta są równe dwóm bokom drugiego trójkąta, to kąty zawarte między równymi bokami są sobie równe.
\end{euclid_theorem}

% Twierdzenie 9.
\begin{euclid_theorem}
    Dany kąt prostokreślny podziel na dwie równe części.
\end{euclid_theorem}

% Twierdzenie 10.
\begin{euclid_theorem}
    Daną linię prostą oznaczoną podzielić na dwie równe części.
\end{euclid_theorem}

% Twierdzenie 11.
\begin{euclid_theorem}
    Z punktu danego na danej linii prostej wyprowadzić linie prostopadłą do danej linii prostej.
\end{euclid_theorem}

% Twierdzenie 12.
\begin{euclid_theorem}
    Z punktu danego leżącego poza linią prostą nieograniczoną, wyprowadzić prostą linię prostopadłą do niej.
\end{euclid_theorem}

% Twierdzenie 13.
\begin{euclid_theorem}
    Jeżeli linia prosta przecinająca drugą prostą tworzy z nią dwa kąty, to są one proste, albo równe dwóm kątom prostym.
\end{euclid_theorem}

% Twierdzenie 14.
\begin{euclid_theorem}
    Jeżeli przy linii prostej i przy punkcie na niej leżącym dwie linie proste nie po jednej stronie położone czynią kąty przyległe równe dwóm kątom prostym, to te linie proste będą w tym samym kierunku.
\end{euclid_theorem}

% Twierdzenie 15.
\begin{euclid_theorem}
    Jeżeli dwie linie proste przecinają się, to utworzone przez nie kąty przeciwległe są sobie równe.
\end{euclid_theorem}

% Twierdzenie 16.
\begin{euclid_theorem}
    W dowolnym trójkącie kąt zewnętrzny powstały przez przedłużenie jednego boku jest większy od każdego z dwóch kątów wewnętrznych przeciwległych jemu.
\end{euclid_theorem}

% Twierdzenie 17.
\begin{euclid_theorem}
    W każdym trójkącie suma dwóch dowolnych kątów jest mniejsza od dwóch kątów prostych.
\end{euclid_theorem}

% Twierdzenie 18.
\begin{euclid_theorem}
    W każdym trójkącie bok większy przeciwległy jest kątowi większemu.
\end{euclid_theorem}

% Twierdzenie 19.
\begin{euclid_theorem}
    W każdym trójkącie kąt większy przeciwległy jest bokowi większemu.
\end{euclid_theorem}

% Twierdzenie 20.
\begin{euclid_theorem}
    W każdym trójkącie suma dwóch dowolnych boków jest większa od boku trzeciego.
\end{euclid_theorem}

% Twierdzenie 21.
\begin{euclid_theorem}
    Jeżeli z końców jednego boku trójkąta poprowadzone będą dwie linie proste wewnątrz trójkąta, aż do zejścia się z sobą, to te dwie linie proste będą mniejsze od dwóch pozostałych boków trójkąta, lecz zawierać jednak 
\end{euclid_theorem}
będą kąt większy od kąta zawartego między pozostałymi bokami trójkąta.
% Twierdzenie 22.
\begin{euclid_theorem}
    Aby z trzech danych linii prostych wykreślić trójkąt, potrzeba aby z tych trzech danych linii prostych suma dwóch którychkolwiek była większa od trzeciej.
\end{euclid_theorem}

% Twierdzenie 23.
\begin{euclid_theorem}
    Na danej linii prostej i punkcie na niej danym wykreślić kąt prostokreślny równy kątowi prostokreślnemu danemu.
\end{euclid_theorem}

% Twierdzenie 24.
\begin{euclid_theorem}
    Jeżeli dwa boki jednego trójkąta, są równe dwóm bokom trójkąta drugiego, z kątów zaś między bokami równymi jeden większy jest od drugiego; to będzie też podstawa jednego trójkąta większa od podstawy drugiego trójkąta.
\end{euclid_theorem}

% Twierdzenie 25.
\begin{euclid_theorem}
    Jeżeli dwa boki jednego trójkąta, są równe dwóm bokom trójkąta drugiego, lecz podstawa jednego trójkąta większa jest od podstawy drugiego trójkąta, to i kąty między bokami równymi zawarte będą jeden większy od drugiego.
\end{euclid_theorem}

% Twierdzenie 26.
\begin{euclid_theorem}
    Jeżeli dwa kąty jednego trójkąta są równe dwóm kątom drugiego trójkąta, i bok jeden przyległy obydwu kątom, albo jednemu w pierwszym trójkącie równa się bokowi jednemu przyległemu obydwu katom, albo jednemu w drugim trójkącie; będą i dwa boki pozostałe równe dwóm bokom pozostałym i kąt trzeci w jednym trójkącie będzie równy katowi trzeciemu w drugim trójkącie.
\end{euclid_theorem}

% Twierdzenie 27.
\begin{euclid_theorem}
    Jeżeli na dwie linie proste, pada linia prosta czyniąca kąty naprzemian równe między sobą, to te dwie linie proste będą równoległe.
\end{euclid_theorem}

% Twierdzenie 28.
\begin{euclid_theorem}
    Jeśli linia prosta opada na dwie linie proste, tworząc kąt zewnętrzny równy wewnętrznemu i przeciwny do kąta na tym samym boku lub suma kątów wewnętrznych na tym samym boku jest równa dwóm kątom prostym, wtedy linie proste są równoległe do siebie.
\end{euclid_theorem}

% Twierdzenie 29.
\begin{euclid_theorem}
    Linia prosta opada na równoległą linie prostą tworząc alternatywne kąty równe sobie, kąt zewnętrzny równy wewnętrznemu i przeciwległy i suma kątów wewnętrznych na tym samym boku jest równa dwóm kątom prostym.
\end{euclid_theorem}

% Twierdzenie 30.
\begin{euclid_theorem}
    Linie proste, które są równoległe do linii prostej są również równoległe do siebie.
\end{euclid_theorem}

% Twierdzenie 31.
\begin{euclid_theorem}
    Poprowadzić przez dany punkt linię prostą równoległą względem danej lini prostej.
\end{euclid_theorem}

% Twierdzenie 32.
\begin{euclid_theorem}
    W jakimkolwiek trójkącie, jeśli jeden z boków jest znany wtedy kąt zewnętrzny jest równy sumie dwóch kątów wewnętrznych i przeciwnych i suma trzech wewnętrznych kątów trójkąta jest równa dwóm kątom prostym.
\end{euclid_theorem}

% Twierdzenie 33.
\begin{euclid_theorem}
    Linie proste, które łączą końce równych i równoległych linii prostych w tym samym kierunku są sobie równe i równoległe.
\end{euclid_theorem}

% Twierdzenie 34.
\begin{euclid_theorem}
    W równoległobokach boki i kąty przeciwne są między sobą równe, a przekątna dzieli je na dwie równe części.
\end{euclid_theorem}

% Twierdzenie 35.
\begin{euclid_theorem}
    Równoległoboki, które są na takiej samej podstawie i są porównywalne są sobie równe.
\end{euclid_theorem}

% Twierdzenie 36.
\begin{euclid_theorem}
    Równoległoboki, które mają równe podstawy i są porównywalne są sobie równe.
\end{euclid_theorem}

% Twierdzenie 37.
\begin{euclid_theorem}
    Trójkąty, które mają takie same podstawy i są porównywalne są sobie równe.
\end{euclid_theorem}

% Twierdzenie 38.
\begin{euclid_theorem}
    Trójkąty, których podstawy są równe i są one porównywalne są sobie równe.
\end{euclid_theorem}

% Twierdzenie 39.
\begin{euclid_theorem}
    Równe trójkąty, które są na takich samych podstawach i mające te same boki również są porównywalne.
\end{euclid_theorem}

% Twierdzenie 40.
\begin{euclid_theorem}
    Równe trójkąty, które mają takie same podstawy i mają te same boki również są porównywalne.
\end{euclid_theorem}

% Twierdzenie 41.
\begin{euclid_theorem}
    Jeśli równoległobok i trójkąt mają tą samą podstawę i są tymi samymi liniami zakończone, to trójkąt jest połową równoległoboku.
\end{euclid_theorem}

% Twierdzenie 42.
\begin{euclid_theorem}
    Skonstruować równoległobok równy danemu trójkątowi o podanym prostoliniowym kącie.
\end{euclid_theorem}

% Twierdzenie 43.
\begin{euclid_theorem}
    W każdym równoległoboku, dopełnienia równoległoboków koło przekątnych położonych są między sobą równe.
\end{euclid_theorem}

% Twierdzenie 44.
\begin{euclid_theorem}
    Na danej linii prostej wykreślić równy danemu równoległobok, którego jeden kąt będzie równy danemu.
\end{euclid_theorem}

% Twierdzenie 45.
\begin{euclid_theorem}
    Wykreślić równy danej figurze prostokreślny równoległobok, którego jeden kąt będzie równy danemu.
\end{euclid_theorem}

% Twierdzenie 46.
\begin{euclid_theorem}
    Na danej linii prostej wykreślić kwadrat.
\end{euclid_theorem}

% Twierdzenie 47.
\begin{euclid_theorem}
    W trójkącie prostokątnym, kwadrat zbudowany na boku przeciwnym kątowi prostemu, równy jest kwadratom zbudowanym na bokach, które kąt prosty zawierają.
\end{euclid_theorem}

% Twierdzenie 48.
\begin{euclid_theorem}
    Jeżeli kwadrat zbudowany na jednym z boków trójkąta, jest równy kwadratom wykreślonym na dwóch pozostałych bokach trójkąta, to kąt zawarty między dwoma pozostałymi bokami będzie prosty.
\end{euclid_theorem}